\shead{Viết số dưới dạng tỉ số phần trăm và ngược lại}
\begin{knowbox}\begin{itemize}[leftmargin=1.4em,itemsep=3pt]
\item Kí hiệu $\%$ có nghĩa là ``phần trăm'': $\ 1\%=\dfrac1{100}=0{,}01$. \ Vì vậy $\ a=a\times100\%$.
\item \textbf{Viết một số dưới dạng tỉ số phần trăm:} nhân số đó với $100$ rồi viết thêm kí hiệu $\%$ vào bên phải.\par
Ví dụ: $\ 3=3\times100\%=300\%;\quad 0{,}425=42{,}5\%;\quad \dfrac34=\dfrac34\times100\%=75\%$.
\item \textbf{Viết tỉ số phần trăm dưới dạng số thập phân:} chia cho $100$ (dời dấu phẩy sang trái hai chữ số).\par
Ví dụ: $\ 36\%=36:100=0{,}36;\quad 4{,}5\%=0{,}045;\quad 120\%=1{,}2$.
\item \textbf{Viết tỉ số phần trăm dưới dạng phân số:} viết thành phân số có mẫu số $100$ rồi rút gọn.\par
Ví dụ: $\ 85\%=\dfrac{85}{100}=\dfrac{17}{20}$.
\end{itemize}
\note{Kí hiệu $\%$ được dùng thay cho phép chia cho $100$.}
\end{knowbox}
\gap
\shead{Các phép tính với tỉ số phần trăm}
\begin{knowbox}\begin{itemize}[leftmargin=1.4em,itemsep=3pt]
\item \textbf{Phép cộng, phép trừ:} cộng, trừ phần số và giữ nguyên kí hiệu $\%$:\par
$a\%+b\%=(a+b)\%;\qquad a\%-b\%=(a-b)\%$.
\item \textbf{Nhân, chia tỉ số phần trăm với một số $k$:} $\ a\%\times k=(a\times k)\%;\quad a\%:k=(a:k)\%$.
\item \textbf{Nhân hai tỉ số phần trăm:} $\ a\%\times b\%=\dfrac a{100}\times\dfrac b{100}$. \ Nên đổi cả hai ra số thập phân rồi tính:\par
$30\%\times40\%=0{,}3\times0{,}4=0{,}12=12\%$ \ (\emph{không phải} $1\,200\%$!).
\item \textbf{Chia hai tỉ số phần trăm:} hai kí hiệu $\%$ ``triệt tiêu'' nhau, kết quả là một số: $\ a\%:b\%=a:b$. \ Ví dụ: $\ 45\%:5\%=9$.
\end{itemize}
\note{Khi chưa chắc chắn, hãy coi mỗi tỉ số phần trăm là một phân số có mẫu số $100$ (hoặc một số thập phân) rồi tính như bình thường.}
\end{knowbox}
\gap
\shead{Biểu thức chứa nhiều loại số}
\begin{knowbox}\begin{itemize}[leftmargin=1.4em,itemsep=3pt]
\item Một biểu thức có thể chứa cùng lúc số thập phân, phân số, hỗn số và tỉ số phần trăm. Trước hết ta đưa tất cả về \textbf{cùng một loại số} (cùng là số thập phân hoặc cùng là phân số), rồi tính theo đúng thứ tự thực hiện phép tính: trong ngoặc trước; nhân, chia trước; cộng, trừ sau.
\item Khi tìm $x$, ta đổi các tỉ số phần trăm ra số thập phân hoặc phân số rồi làm như các bài tìm thành phần chưa biết đã học.
\end{itemize}
\end{knowbox}
\gap
\shead{Tỉ số phần trăm trong bài toán có lời văn}
\begin{knowbox}\begin{itemize}[leftmargin=1.4em,itemsep=3pt]
\item ``$A$ bằng $a\%$ của $B$'' nghĩa là $A$ bằng $\dfrac a{100}$ của $B$. \ Rút gọn phân số để thấy số phần:\par
``$A$ bằng $40\%$ của $B$'' $\ \Rightarrow\ A=\dfrac25\ B$: \ $A$ gồm $2$ phần, $B$ gồm $5$ phần bằng nhau.
\item Sau đó vẽ sơ đồ đoạn thẳng và dùng các cách giải quen thuộc: \textbf{tổng -- tỉ}, \textbf{hiệu -- tỉ}.
\item Cả một đại lượng ứng với $100\%$. \ Nếu biết các phần dưới dạng tỉ số phần trăm thì phần còn lại bằng $100\%$ trừ đi các tỉ số đã biết.
\end{itemize}
\end{knowbox}
\gap
\begin{exbox}{\ExWord\ 1}
\pa{a} Viết các số sau dưới dạng tỉ số phần trăm: $\ 0{,}36;\quad \dfrac7{20};\quad 3;\quad 1{,}075$.\par
\pa{b} Viết các tỉ số phần trăm sau dưới dạng số thập phân và dưới dạng phân số tối giản: $\ 45\%;\quad 6{,}5\%;\quad 160\%$.
\guide
\pa{a} $0{,}36=0{,}36\times100\%=36\%$.\par
$\dfrac7{20}=\dfrac{7\times5}{20\times5}=\dfrac{35}{100}=35\%$ \ \ hoặc \ \ $\dfrac7{20}=7:20=0{,}35=35\%$.\par
$3=3\times100\%=300\%$; \qquad $1{,}075=107{,}5\%$.\par
\pa{b} $45\%=45:100=0{,}45$; \quad $45\%=\dfrac{45}{100}=\dfrac9{20}$.\par
$6{,}5\%=6{,}5:100=0{,}065$; \quad $6{,}5\%=\dfrac{65}{1\,000}=\dfrac{13}{200}$.\par
$160\%=1{,}6$; \quad $160\%=\dfrac{160}{100}=\dfrac85$.
\end{exbox}
\gap
\begin{exbox}{\ExWord\ 2}
Tính.
\begin{cols3}
\pa{a} $37\%+48\%$ & \pa{b} $126\%-58\%$ & \pa{c} $15{,}5\%\times6$\\
\pa{d} $40\%\times15\%$ & \pa{e} $72\%:9\%$ & \pa{g} $72\%:9$
\end{cols3}
\guide
\pa{a} \emph{Cách 1.} $37\%+48\%=(37+48)\%=85\%$.\par
\emph{Cách 2.} $37\%+48\%=0{,}37+0{,}48=0{,}85=85\%$.\par
\pa{b} $126\%-58\%=(126-58)\%=68\%$.\par
\pa{c} $15{,}5\%\times6=(15{,}5\times6)\%=93\%$.\par
\pa{d} $40\%\times15\%=0{,}4\times0{,}15=0{,}06=6\%$.\par
\pa{e} $72\%:9\%=0{,}72:0{,}09=8$.\par
\pa{g} $72\%:9=(72:9)\%=8\%$.
\note{Chú ý: $40\%\times15\%$ \emph{không} bằng $600\%$; còn $72\%:9\%$ là số $8$ chứ không phải $8\%$.}
\end{exbox}
\gap
\begin{exbox}{\ExWord\ 3}
Tính: $\ A=4{,}5-1\dfrac34+63\%:18\%$.
\guide
\emph{Cách 1 (đưa về số thập phân).}\par
$A=4{,}5-1{,}75+0{,}63:0{,}18=4{,}5-1{,}75+3{,}5=2{,}75+3{,}5=6{,}25$.\par
\emph{Cách 2 (đưa về phân số).}\par
$A=\dfrac92-\dfrac74+\dfrac{63}{100}:\dfrac{18}{100}=\dfrac92-\dfrac74+\dfrac72=\dfrac{18}4-\dfrac74+\dfrac{14}4=\dfrac{25}4=6\dfrac14$.
\note{Viết kết quả $A=6{,}25$ hay $A=6\dfrac14$ đều đúng.}
\end{exbox}
\gap
\begin{exbox}{\ExWord\ 4}
Tìm $x$, biết: $\ x:25\%+1{,}6=9:0{,}5$.
\guide
Tính vế phải: $9:0{,}5=18$.\par
Ta có $x:25\%+1{,}6=18$, \ nên $x:25\%=18-1{,}6=16{,}4$.\par
Suy ra $x=16{,}4\times25\%=16{,}4\times0{,}25=4{,}1$.\par
\emph{Thử lại:} $4{,}1:0{,}25+1{,}6=16{,}4+1{,}6=18$.
\dapso{$x=4{,}1$.}
\end{exbox}
\gap
\begin{exbox}{\ExWord\ 5}
Lan và Hoa có tất cả $72$ nhãn dán. Số nhãn dán của Hoa bằng $80\%$ số nhãn dán của Lan. Hỏi mỗi bạn có bao nhiêu nhãn dán?
\guide
$80\%=\dfrac{80}{100}=\dfrac45$, nên số nhãn dán của Hoa bằng $\dfrac45$ số nhãn dán của Lan.
\begin{barfigC}
\rowlab{-0.15}{0.6}{Lan:}
\foreach \i in {0,...,4}{\blk{\i*0.8}{\i*0.8+0.8}{0.6}{0.42}{tmBlue}{}}
\rowlab{-0.15}{0}{Hoa:}
\foreach \i in {0,...,3}{\blk{\i*0.8}{\i*0.8+0.8}{0}{0.42}{tmGreen}{}}
\vbrace{4.1}{0.81}{-0.21}{$72$ nhãn dán}
\end{barfigC}
Tổng số phần bằng nhau: $5+4=9$ (phần).\par
Giá trị một phần: $72:9=8$ (nhãn dán).\par
Lan có: $8\times5=40$ (nhãn dán); \ Hoa có: $72-40=32$ (nhãn dán).
\dapso{Lan $40$ nhãn dán, Hoa $32$ nhãn dán.}
\end{exbox}
\dividerbar{Luyện tập}
\begin{probbox}{1}{\relax}
Viết các số sau dưới dạng tỉ số phần trăm.
\begin{cols3}
\pa{a} $0{,}47$ & \pa{b} $1{,}36$ & \pa{c} $0{,}209$\\
\pa{d} $0{,}8$ & \pa{e} $2{,}5$ & \pa{g} $0{,}035$
\end{cols3}
\end{probbox}
\probsol{
Nhân với $100$ rồi viết thêm kí hiệu $\%$:\par
\pa{a} $0{,}47=47\%$; \quad \pa{b} $1{,}36=136\%$; \quad \pa{c} $0{,}209=20{,}9\%$;\par
\pa{d} $0{,}8=80\%$; \quad \pa{e} $2{,}5=250\%$; \quad \pa{g} $0{,}035=3{,}5\%$.
\dapso{$47\%$; $136\%$; $20{,}9\%$; $80\%$; $250\%$; $3{,}5\%$.}}
\gap
\begin{probbox}{2}{\relax}
Viết các phân số sau dưới dạng tỉ số phần trăm.
\begin{cols3}
\pa{a} $\dfrac35$ & \pa{b} $\dfrac9{20}$ & \pa{c} $\dfrac{21}{30}$\\
\pa{d} $\dfrac78$ & \pa{e} $\dfrac{13}4$ & \pa{g} $\dfrac{33}{40}$
\end{cols3}
\end{probbox}
\probsol{
\pa{a} $\dfrac35=\dfrac{60}{100}=60\%$. \quad \pa{b} $\dfrac9{20}=\dfrac{45}{100}=45\%$. \quad \pa{c} $\dfrac{21}{30}=\dfrac7{10}=\dfrac{70}{100}=70\%$.\par
\pa{d} $\dfrac78=7:8=0{,}875=87{,}5\%$. \quad \pa{e} $\dfrac{13}4=\dfrac{325}{100}=325\%$.\par
\pa{g} $\dfrac{33}{40}=33:40=0{,}825=82{,}5\%$.
\dapso{$60\%$; $45\%$; $70\%$; $87{,}5\%$; $325\%$; $82{,}5\%$.}}
\gap
\begin{probbox}{3}{\relax}
Viết các tỉ số phần trăm sau dưới dạng số thập phân.
\begin{cols3}
\pa{a} $38\%$ & \pa{b} $4\%$ & \pa{c} $215\%$\\
\pa{d} $9{,}6\%$ & \pa{e} $300\%$ & \pa{g} $0{,}5\%$
\end{cols3}
\end{probbox}
\probsol{
Chia cho $100$:\par
\pa{a} $38\%=0{,}38$; \quad \pa{b} $4\%=0{,}04$; \quad \pa{c} $215\%=2{,}15$;\par
\pa{d} $9{,}6\%=0{,}096$; \quad \pa{e} $300\%=3$; \quad \pa{g} $0{,}5\%=0{,}005$.
\dapso{$0{,}38$; $0{,}04$; $2{,}15$; $0{,}096$; $3$; $0{,}005$.}}
\gap
\begin{probbox}{4}{\relax}
Viết các tỉ số phần trăm sau dưới dạng phân số tối giản.
\begin{cols3}
\pa{a} $40\%$ & \pa{b} $85\%$ & \pa{c} $64\%$\\
\pa{d} $150\%$ & \pa{e} $2{,}5\%$ & \pa{g} $62{,}5\%$
\end{cols3}
\end{probbox}
\probsol{
\pa{a} $40\%=\dfrac{40}{100}=\dfrac25$. \quad \pa{b} $85\%=\dfrac{85}{100}=\dfrac{17}{20}$. \quad \pa{c} $64\%=\dfrac{64}{100}=\dfrac{16}{25}$.\par
\pa{d} $150\%=\dfrac{150}{100}=\dfrac32$. \quad \pa{e} $2{,}5\%=\dfrac{25}{1\,000}=\dfrac1{40}$. \quad \pa{g} $62{,}5\%=\dfrac{625}{1\,000}=\dfrac58$.
\dapso{$\dfrac25$; $\dfrac{17}{20}$; $\dfrac{16}{25}$; $\dfrac32$; $\dfrac1{40}$; $\dfrac58$.}}
\gap
\begin{probbox}{5}{\relax}
Điền dấu $>$, $<$, $=$ thích hợp vào chỗ chấm.
\begin{cols2}
\pa{a} $0{,}7\ \ \dots\ \ 69\%$ & \pa{b} $\dfrac38\ \ \dots\ \ 37{,}5\%$\\
\pa{c} $125\%\ \ \dots\ \ 1\dfrac15$ & \pa{d} $\dfrac23\ \ \dots\ \ 66\%$
\end{cols2}
\end{probbox}
\probsol{
\pa{a} $0{,}7=70\%$ mà $70\%>69\%$, nên $0{,}7>69\%$.\par
\pa{b} $\dfrac38=3:8=0{,}375=37{,}5\%$, nên $\dfrac38=37{,}5\%$.\par
\pa{c} $125\%=1{,}25$ và $1\dfrac15=1{,}2$; \ $1{,}25>1{,}2$, nên $125\%>1\dfrac15$.\par
\pa{d} $66\%=\dfrac{66}{100}=\dfrac{33}{50}$, \ mà $\dfrac23=\dfrac{100}{150}>\dfrac{99}{150}=\dfrac{33}{50}$, nên $\dfrac23>66\%$.
\dapso{\pa{a} $>$;\ \ \pa{b} $=$;\ \ \pa{c} $>$;\ \ \pa{d} $>$.}}
\gap
\begin{probbox}{6}{\relax}
Sắp xếp các số sau theo thứ tự từ bé đến lớn: $\ 0{,}43;\ \ \dfrac25;\ \ 41{,}5\%;\ \ \dfrac9{20};\ \ 39\%$.
\end{probbox}
\probsol{
Đổi tất cả ra số thập phân:\par
$0{,}43$; \quad $\dfrac25=0{,}4$; \quad $41{,}5\%=0{,}415$; \quad $\dfrac9{20}=0{,}45$; \quad $39\%=0{,}39$.\par
Ta có: $0{,}39<0{,}4<0{,}415<0{,}43<0{,}45$.
\dapso{$39\%;\ \dfrac25;\ 41{,}5\%;\ 0{,}43;\ \dfrac9{20}$.}}
\gap
\begin{probbox}{7}{\relax}
Tính.
\begin{cols2}
\pa{a} $28\%+56\%$ & \pa{b} $143\%-67\%$\\
\pa{c} $12{,}5\%+7{,}25\%+30{,}25\%$ & \pa{d} $90\%-45\%+15\%$
\end{cols2}
\end{probbox}
\probsol{
\pa{a} $28\%+56\%=(28+56)\%=84\%$.\par
\pa{b} $143\%-67\%=(143-67)\%=76\%$.\par
\pa{c} $12{,}5\%+7{,}25\%+30{,}25\%=(12{,}5+7{,}25+30{,}25)\%=50\%$.\par
\pa{d} $90\%-45\%+15\%=(90-45+15)\%=60\%$.
\dapso{$84\%$; $76\%$; $50\%$; $60\%$.}}
\gap
\begin{probbox}{8}{\relax}
Tính.
\begin{cols2}
\pa{a} $23{,}4\%\times5$ & \pa{b} $252\%:12$\\
\pa{c} $6{,}8\%\times15$ & \pa{d} $99\%:4{,}5$
\end{cols2}
\end{probbox}
\probsol{
\pa{a} $23{,}4\%\times5=(23{,}4\times5)\%=117\%$.\par
\pa{b} $252\%:12=(252:12)\%=21\%$.\par
\pa{c} $6{,}8\%\times15=(6{,}8\times15)\%=102\%$.\par
\pa{d} $99\%:4{,}5=(99:4{,}5)\%=22\%$.
\dapso{$117\%$; $21\%$; $102\%$; $22\%$.}}
\gap
\begin{probbox}{9}{\relax}
Thực hiện phép tính rồi viết kết quả dưới dạng số thập phân.
\begin{cols3}
\pa{a} $60\%\times35\%$ & \pa{b} $140\%\times25\%$ & \pa{c} $3{,}5\%\times80\%$\\
\pa{d} $84\%:12\%$ & \pa{e} $45\%:0{,}5\%$ & \pa{g} $7{,}2\%:16\%$
\end{cols3}
\end{probbox}
\probsol{
\pa{a} $60\%\times35\%=0{,}6\times0{,}35=0{,}21$.\par
\pa{b} $140\%\times25\%=1{,}4\times0{,}25=0{,}35$.\par
\pa{c} $3{,}5\%\times80\%=0{,}035\times0{,}8=0{,}028$.\par
\pa{d} $84\%:12\%=84:12=7$.\par
\pa{e} $45\%:0{,}5\%=45:0{,}5=90$.\par
\pa{g} $7{,}2\%:16\%=7{,}2:16=0{,}45$.
\dapso{$0{,}21$; $0{,}35$; $0{,}028$; $7$; $90$; $0{,}45$.}}
\gap
\begin{probbox}{10}{\relax}
Tính.
\begin{cols2}
\pa{a} $2{,}4+35\%-\dfrac34$ & \pa{b} $1\dfrac12\times40\%+2{,}8:70\%$
\end{cols2}
\end{probbox}
\probsol{
\pa{a} $2{,}4+35\%-\dfrac34=2{,}4+0{,}35-0{,}75=2{,}75-0{,}75=2$.\par
\pa{b} $1\dfrac12\times40\%+2{,}8:70\%=1{,}5\times0{,}4+2{,}8:0{,}7=0{,}6+4=4{,}6$.
\dapso{\pa{a} $2$;\ \ \pa{b} $4{,}6$.}}
\gap
\begin{probbox}{11}{\relax}
Tính.
\begin{cols2}
\pa{a} $25\%\times8+13{,}5:4{,}5-2\dfrac25$ & \pa{b} $4{,}8-1\dfrac14+36\%:144\%$
\end{cols2}
\end{probbox}
\probsol{
\pa{a} $25\%\times8+13{,}5:4{,}5-2\dfrac25=0{,}25\times8+3-2{,}4=2+3-2{,}4=2{,}6$.\par
\pa{b} $4{,}8-1\dfrac14+36\%:144\%=4{,}8-1{,}25+36:144=4{,}8-1{,}25+0{,}25=3{,}8$.
\dapso{\pa{a} $2{,}6$;\ \ \pa{b} $3{,}8$.}}
\gap
\begin{probbox}{12}{\relax}
Tính bằng cách thuận tiện.
\begin{cols3}
\pa{a} $37\%\times2{,}5+63\%\times2{,}5$ & \pa{b} $125\%\times4{,}7\times80\%$ & \pa{c} $6{,}4\times45\%+6{,}4\times55\%$
\end{cols3}
\end{probbox}
\probsol{
\pa{a} $37\%\times2{,}5+63\%\times2{,}5=(37\%+63\%)\times2{,}5=100\%\times2{,}5=1\times2{,}5=2{,}5$.\par
\pa{b} $125\%\times4{,}7\times80\%=(1{,}25\times0{,}8)\times4{,}7=1\times4{,}7=4{,}7$.\par
\pa{c} $6{,}4\times45\%+6{,}4\times55\%=6{,}4\times(45\%+55\%)=6{,}4\times1=6{,}4$.
\dapso{$2{,}5$; $4{,}7$; $6{,}4$.}}
\gap
\begin{probbox}{13}{\relax}
Tìm $x$, biết:
\begin{cols2}
\pa{a} $20\%\times x=1{,}35+0{,}25$ & \pa{b} $x:40\%-2{,}5=7{,}5:1{,}5$
\end{cols2}
\end{probbox}
\probsol{
\pa{a} $0{,}2\times x=1{,}6$, \ nên $x=1{,}6:0{,}2=8$.\par
\pa{b} $x:40\%-2{,}5=5$, \ nên $x:40\%=7{,}5$ \ và \ $x=7{,}5\times0{,}4=3$.
\dapso{\pa{a} $x=8$;\ \ \pa{b} $x=3$.}}
\gap
\begin{probbox}{14}{\relax}
Tìm $x$, biết:
\begin{cols2}
\pa{a} $150\%\times x=4{,}2-1{,}8$ & \pa{b} $x:25\%+0{,}75=0{,}75\times21$
\end{cols2}
\end{probbox}
\probsol{
\pa{a} $1{,}5\times x=2{,}4$, \ nên $x=2{,}4:1{,}5=1{,}6$.\par
\pa{b} $x:25\%+0{,}75=15{,}75$, \ nên $x:25\%=15$ \ và \ $x=15\times0{,}25=3{,}75$.
\dapso{\pa{a} $x=1{,}6$;\ \ \pa{b} $x=3{,}75$.}}
\gap
\begin{probbox}{15}{\relax}
Tìm $x$, biết:
\begin{cols2}
\pa{a} $x+35\%\times x=5{,}4$ & \pa{b} $x-12\%\times x=22$
\end{cols2}
\end{probbox}
\probsol{
\pa{a} $x=100\%\times x$, nên $x+35\%\times x=135\%\times x=1{,}35\times x$.\par
$1{,}35\times x=5{,}4$, \ nên $x=5{,}4:1{,}35=4$.\par
\pa{b} $x-12\%\times x=88\%\times x=0{,}88\times x=22$, \ nên $x=22:0{,}88=25$.
\dapso{\pa{a} $x=4$;\ \ \pa{b} $x=25$.}}
\gap
\begin{probbox}{16}{\relax}
Một vòi nước chảy vào một bể chưa có nước. Giờ thứ nhất vòi chảy được $30\%$ bể, giờ thứ hai chảy được nhiều hơn giờ thứ nhất $12\%$ bể. Hỏi sau hai giờ, còn bao nhiêu phần trăm bể chưa có nước?
\end{probbox}
\probsol{
Giờ thứ hai vòi chảy được: $30\%+12\%=42\%$ (bể).\par
Cả hai giờ chảy được: $30\%+42\%=72\%$ (bể).\par
Phần bể chưa có nước: $100\%-72\%=28\%$ (bể).
\dapso{$28\%$ bể.}}
\gap
\begin{probbox}{17}{\relax}
Thư viện của một trường có ba loại sách: truyện, sách khoa học và sách giáo khoa. Số truyện chiếm $45\%$ tổng số sách; số sách khoa học bằng $\dfrac23$ số truyện. Hỏi sách giáo khoa chiếm bao nhiêu phần trăm tổng số sách?
\end{probbox}
\probsol{
Sách khoa học chiếm: $45\%\times\dfrac23=30\%$ (tổng số sách).\par
Sách giáo khoa chiếm: $100\%-45\%-30\%=25\%$ (tổng số sách).
\dapso{$25\%$.}}
\gap
\begin{probbox}{18}{\relax}
Năm nay tuổi Khoa bằng $25\%$ tuổi bố, tuổi ông bằng $175\%$ tuổi bố. Tổng số tuổi của Khoa và bố là $50$ tuổi. Tính tuổi của mỗi người hiện nay.
\end{probbox}
\probsol{
$25\%=\dfrac14$: tuổi Khoa là $1$ phần, tuổi bố là $4$ phần.\par
Tổng số phần bằng nhau: $1+4=5$ (phần); \ giá trị một phần: $50:5=10$ (tuổi).\par
Tuổi Khoa: $10$ tuổi; \ tuổi bố: $10\times4=40$ (tuổi).\par
Tuổi ông: $40\times175\%=40\times1{,}75=70$ (tuổi).
\dapso{Khoa $10$ tuổi, bố $40$ tuổi, ông $70$ tuổi.}}
\gap
\begin{probbox}{19}{\relax}
Một hình chữ nhật có chu vi $32$ dm, chiều rộng bằng $60\%$ chiều dài. Tính diện tích hình chữ nhật đó.
\end{probbox}
\probsol{
Nửa chu vi (tổng chiều dài và chiều rộng): $32:2=16$ (dm).\par
$60\%=\dfrac35$: chiều rộng là $3$ phần, chiều dài là $5$ phần.\par
Giá trị một phần: $16:(3+5)=2$ (dm). \ Chiều dài: $2\times5=10$ (dm); \ chiều rộng: $2\times3=6$ (dm).\par
Diện tích: $10\times6=60$ (dm$^2$).
\dapso{$60$ dm$^2$.}}
\gap
\begin{probbox}{20}{\relax}
Một trang trại nuôi thỏ và gà, tất cả $150$ con. Sau khi bán $10$ con thỏ và $20$ con gà thì số thỏ còn lại bằng $50\%$ số gà còn lại. Hỏi lúc đầu trang trại có bao nhiêu con thỏ, bao nhiêu con gà?
\end{probbox}
\probsol{
Số con còn lại: $150-10-20=120$ (con).\par
$50\%=\dfrac12$: số thỏ còn lại là $1$ phần, số gà còn lại là $2$ phần.\par
Giá trị một phần: $120:3=40$ (con). \ Thỏ còn: $40$ con; \ gà còn: $80$ con.\par
Lúc đầu: thỏ có $40+10=50$ (con); \ gà có $80+20=100$ (con).
\dapso{$50$ con thỏ và $100$ con gà.}}
\gap
\begin{probbox}{21}{\relax}
Hai ngăn sách có tất cả $72$ quyển. Nếu chuyển $20\%$ số sách ở ngăn thứ nhất sang ngăn thứ hai thì số sách ở hai ngăn bằng nhau. Hỏi mỗi ngăn có bao nhiêu quyển sách?
\end{probbox}
\probsol{
Sau khi chuyển, mỗi ngăn có: $72:2=36$ (quyển).\par
Ngăn thứ nhất còn lại $100\%-20\%=80\%$ số sách của nó, nên $80\%$ số sách ngăn thứ nhất là $36$ quyển.\par
Ngăn thứ nhất có: $36:80\%=36:0{,}8=45$ (quyển).\par
Ngăn thứ hai có: $72-45=27$ (quyển).\par
\emph{Thử lại:} $20\%$ của $45$ là $9$; \ $45-9=36$ và $27+9=36$.
\dapso{$45$ quyển và $27$ quyển.}}
\gap
\begin{probbox}{22}{\relax}
Hai kho A và B chứa gạo. Nếu chuyển $2$ tấn gạo từ kho A sang kho B thì số gạo ở hai kho bằng nhau. Nếu chuyển $1$ tấn gạo từ kho B sang kho A thì số gạo ở kho A bằng $300\%$ số gạo ở kho B. Hỏi lúc đầu mỗi kho có bao nhiêu tấn gạo?
\end{probbox}
\probsol{
Chuyển $2$ tấn từ A sang B thì hai kho bằng nhau, nên kho A nhiều hơn kho B: $2\times2=4$ (tấn).\par
Nếu chuyển $1$ tấn từ B sang A thì hiệu hai kho là: $4+1+1=6$ (tấn).\par
Khi đó A gấp $300\%=3$ lần B: B là $1$ phần, A là $3$ phần, hiệu là $2$ phần.\par
Giá trị một phần: $6:2=3$ (tấn). \ Khi đó kho B có $3$ tấn, nên lúc đầu kho B có: $3+1=4$ (tấn).\par
Lúc đầu kho A có: $4+4=8$ (tấn).
\dapso{Kho A $8$ tấn, kho B $4$ tấn.}}
\gap
\begin{probbox}{23}{\relax}
Đầu năm học, số học sinh nam của lớp 5A bằng $80\%$ số học sinh nữ. Đến giữa năm, có thêm $2$ học sinh nam chuyển đến nên số học sinh nam lúc này bằng $90\%$ số học sinh nữ. Hỏi đầu năm lớp 5A có bao nhiêu học sinh?
\end{probbox}
\probsol{
Số học sinh nữ không thay đổi.\par
$2$ học sinh nam mới chuyển đến ứng với: $90\%-80\%=10\%$ (số học sinh nữ).\par
Số học sinh nữ: $2:10\%=2:0{,}1=20$ (học sinh).\par
Số học sinh nam đầu năm: $20\times80\%=16$ (học sinh).\par
Đầu năm lớp 5A có: $20+16=36$ (học sinh).
\dapso{$36$ học sinh.}}
\gap
\begin{probbox}{24}{\relax}
Một cửa hàng có một số vải. Thứ Hai cửa hàng bán được $18{,}5\%$ số vải, thứ Ba bán được $26{,}5\%$ số vải, thứ Tư bán được bằng cả hai ngày thứ Hai và thứ Ba. Hỏi cửa hàng còn lại bao nhiêu phần trăm số vải?
\end{probbox}
\probsol{
Hai ngày thứ Hai và thứ Ba bán được: $18{,}5\%+26{,}5\%=45\%$ (số vải).\par
Thứ Tư bán được: $45\%$ (số vải).\par
Còn lại: $100\%-45\%-45\%=10\%$ (số vải).
\dapso{$10\%$.}}
\gap
\begin{probbox}{25}{\relax}
Tìm $x$, biết:
\begin{cols2}
\pa{a} $75\%\times x+x\times\dfrac14=3{,}6$ & \pa{b} $x:20\%-x=12$
\end{cols2}
\end{probbox}
\probsol{
\pa{a} $\dfrac14=25\%$, nên $75\%\times x+25\%\times x=100\%\times x=x$. \ Vậy $x=3{,}6$.\par
\pa{b} $x:20\%=x:0{,}2=x\times5$, nên $x\times5-x=x\times4=12$ \ và \ $x=3$.
\dapso{\pa{a} $x=3{,}6$;\ \ \pa{b} $x=3$.}}
\gap
\begin{probbox}{26}{\relax}
Tính: $\ (45\%+0{,}55)\times3{,}7-75\%\times4$.
\end{probbox}
\probsol{
$45\%+0{,}55=0{,}45+0{,}55=1$, \ và $75\%\times4=0{,}75\times4=3$.\par
Vậy biểu thức có giá trị: $1\times3{,}7-3=0{,}7$.
\dapso{$0{,}7$.}}
\gap
\begin{probbox}{27}{\relax}
Ở một trường, $60\%$ số học sinh là nữ, và $25\%$ số học sinh nữ tham gia đội hợp xướng của trường. Hỏi số học sinh nữ trong đội hợp xướng chiếm bao nhiêu phần trăm số học sinh toàn trường?
\end{probbox}
\probsol{
Số học sinh nữ trong đội hợp xướng bằng $25\%$ của $60\%$ số học sinh toàn trường:\par
$60\%\times25\%=0{,}6\times0{,}25=0{,}15=15\%$.
\dapso{$15\%$ số học sinh toàn trường.}}
\gap
\begin{probbox}{28}{\relax}
Tổng của hai số là $2{,}7$. Số thứ nhất bằng $125\%$ số thứ hai. Tìm hai số đó.
\end{probbox}
\probsol{
$125\%=\dfrac{125}{100}=\dfrac54$: số thứ nhất là $5$ phần, số thứ hai là $4$ phần.\par
Giá trị một phần: $2{,}7:(5+4)=0{,}3$.\par
Số thứ nhất: $0{,}3\times5=1{,}5$; \ số thứ hai: $0{,}3\times4=1{,}2$.
\dapso{$1{,}5$ và $1{,}2$.}}
\gap
\begin{probbox}{29}{\relax}
Hiệu của hai số là $3{,}6$. Số bé bằng $40\%$ số lớn. Tìm hai số đó.
\end{probbox}
\probsol{
$40\%=\dfrac25$: số bé là $2$ phần, số lớn là $5$ phần.\par
Hiệu số phần bằng nhau: $5-2=3$ (phần); giá trị một phần: $3{,}6:3=1{,}2$.\par
Số bé: $1{,}2\times2=2{,}4$; \ số lớn: $1{,}2\times5=6$.
\dapso{$2{,}4$ và $6$.}}
\gap
\begin{probbox}{30}{\relax}
Mẹ cắt đi $35\%$ một cuộn ruy băng, rồi cắt tiếp $40\%$ phần còn lại. \ \pa{a} Phần ruy băng còn lại bằng bao nhiêu phần trăm cả cuộn? \ \pa{b} Biết phần còn lại dài $7{,}8$ m. Hỏi lúc đầu cuộn ruy băng dài bao nhiêu mét?
\end{probbox}
\probsol{
\pa{a} Sau lần cắt thứ nhất còn lại: $100\%-35\%=65\%$ (cuộn).\par
Lần cắt thứ hai để lại $100\%-40\%=60\%$ phần đó: $\ 65\%\times60\%=0{,}65\times0{,}6=0{,}39=39\%$ (cuộn).\par
\pa{b} $39\%$ cuộn ruy băng dài $7{,}8$ m, nên cả cuộn dài: $7{,}8:0{,}39=20$ (m).
\dapso{\pa{a} $39\%$;\ \ \pa{b} $20$ m.}}
\dividerbar{Thử thách -- Nâng cao}
\begin{probbox}{31}{\stars{1}}
Ở một bể bơi, số trẻ em ngồi trên bờ bằng $25\%$ số trẻ em dưới nước. Sau khi có $4$ em trên bờ nhảy xuống nước thì số trẻ em trên bờ bằng $12{,}5\%$ số trẻ em dưới nước. Hỏi có tất cả bao nhiêu trẻ em?
\end{probbox}
\probsol{
$25\%=\dfrac14$: lúc đầu số em trên bờ bằng $\dfrac1{1+4}=\dfrac15$ tổng số trẻ em.\par
$12{,}5\%=\dfrac18$: lúc sau số em trên bờ bằng $\dfrac1{1+8}=\dfrac19$ tổng số trẻ em.\par
Tổng số trẻ em không đổi. $4$ em nhảy xuống nước ứng với: $\dfrac15-\dfrac19=\dfrac4{45}$ (tổng số trẻ em).\par
Tổng số trẻ em: $4:\dfrac4{45}=45$ (em).\par
\emph{Thử lại:} lúc đầu $9$ em trên bờ, $36$ em dưới nước ($9$ bằng $25\%$ của $36$); lúc sau $5$ em và $40$ em ($5$ bằng $12{,}5\%$ của $40$).
\dapso{$45$ trẻ em.}}
\gap
\begin{probbox}{32}{\stars{1}}
Số tem của An bằng $75\%$ số tem của Bình. Nếu mỗi bạn có thêm $10$ con tem thì số tem của An bằng $80\%$ số tem của Bình. Hỏi hiện nay mỗi bạn có bao nhiêu con tem?
\end{probbox}
\probsol{
$75\%=\dfrac34$: An có $3$ phần, Bình có $4$ phần, nên Bình hơn An $1$ phần.\par
Khi mỗi bạn có thêm $10$ con tem, hiệu số tem vẫn là $1$ phần đó.\par
$80\%=\dfrac45$: lúc sau An có $4$ phần mới, Bình có $5$ phần mới, hiệu là $1$ phần mới.\par
Vậy $1$ phần cũ $=1$ phần mới, và số tem của Bình tăng từ $4$ phần lên $5$ phần: $10$ con tem ứng với $1$ phần.\par
Hiện nay: An có $3\times10=30$ (con tem), Bình có $4\times10=40$ (con tem).\par
\emph{Thử lại:} $40:50=0{,}8=80\%$.
\dapso{An $30$ con tem, Bình $40$ con tem.}}
\gap
\begin{probbox}{33}{\stars{1}}
Giá một món đồ chơi được tăng thêm $20\%$. Sau đó, giá mới lại được giảm đi $20\%$. So với giá ban đầu, giá cuối cùng tăng hay giảm, và tăng (giảm) bao nhiêu phần trăm?
\end{probbox}
\probsol{
Sau khi tăng, giá bằng $120\%$ giá ban đầu.\par
Sau khi giảm, giá bằng $80\%$ của giá đó: $\ 120\%\times80\%=1{,}2\times0{,}8=0{,}96=96\%$ giá ban đầu.\par
Vậy giá đã giảm: $100\%-96\%=4\%$.
\dapso{giảm $4\%$.}}
\gap
\begin{probbox}{34}{\stars{1}}
Một hình chữ nhật được tăng chiều dài thêm $25\%$ và giảm chiều rộng đi $20\%$. Hỏi diện tích hình chữ nhật thay đổi thế nào?
\end{probbox}
\probsol{
Chiều dài mới bằng $125\%$ chiều dài cũ; chiều rộng mới bằng $80\%$ chiều rộng cũ.\par
Diện tích mới bằng: $125\%\times80\%=1{,}25\times0{,}8=1=100\%$ diện tích cũ.
\dapso{diện tích không thay đổi.}}
\gap
\begin{probbox}{35}{\stars{1}}
Tổng của ba số là $108$. Số thứ nhất bằng $50\%$ số thứ hai, số thứ ba bằng $150\%$ số thứ hai. Tìm ba số đó.
\end{probbox}
\probsol{
$50\%=\dfrac12$ và $150\%=\dfrac32$. \ Coi số thứ hai là $2$ phần thì số thứ nhất là $1$ phần, số thứ ba là $3$ phần.\par
Giá trị một phần: $108:(1+2+3)=18$.\par
Ba số là: $18$; \ $18\times2=36$; \ $18\times3=54$.
\dapso{$18;\ 36;\ 54$.}}
\gap
\begin{probbox}{36}{\stars{1}}
Tìm $x$, biết: $\ (x-15\%\times x):1{,}7=30\%$.
\end{probbox}
\probsol{
$x-15\%\times x=85\%\times x=0{,}85\times x$.\par
$0{,}85\times x=30\%\times1{,}7=0{,}3\times1{,}7=0{,}51$.\par
$x=0{,}51:0{,}85=0{,}6$.
\dapso{$x=0{,}6$.}}
\gap
\begin{probbox}{37}{\stars{1}}
Tính: $\ A=1\%+2\%+3\%+\dots+98\%+99\%$. \ Viết kết quả dưới dạng số thập phân.
\end{probbox}
\probsol{
$A=(1+2+3+\dots+99)\%$.\par
$1+2+\dots+99=(1+99)\times99:2=4\,950$.\par
$A=4\,950\%=49{,}5$.
\dapso{$49{,}5$.}}
\gap
\begin{probbox}{38}{\stars{1}}
Tính bằng cách thuận tiện: $\ 12{,}5\%\times0{,}25\times7{,}9\times8\times400\%$.
\end{probbox}
\probsol{
Nhóm các thừa số: $(12{,}5\%\times8)\times(0{,}25\times400\%)\times7{,}9$.\par
$12{,}5\%\times8=0{,}125\times8=1$; \quad $0{,}25\times400\%=0{,}25\times4=1$.\par
Kết quả: $1\times1\times7{,}9=7{,}9$.
\dapso{$7{,}9$.}}
\gap
\begin{probbox}{39}{\stars{1}}
Một hộp có bi đỏ và bi xanh, trong đó số bi đỏ chiếm $60\%$ tổng số bi. Sau khi bỏ thêm $20$ viên bi xanh vào hộp thì số bi đỏ chiếm $40\%$ tổng số bi. Hỏi trong hộp có bao nhiêu viên bi đỏ?
\end{probbox}
\probsol{
Số bi đỏ không thay đổi.\par
Lúc đầu: bi đỏ chiếm $60\%$, bi xanh chiếm $40\%$, nên số bi xanh bằng $\dfrac{40}{60}=\dfrac23$ số bi đỏ.\par
Lúc sau: bi đỏ chiếm $40\%$, bi xanh chiếm $60\%$, nên số bi xanh bằng $\dfrac{60}{40}=\dfrac32$ số bi đỏ.\par
$20$ viên bi xanh bỏ thêm ứng với: $\dfrac32-\dfrac23=\dfrac56$ (số bi đỏ).\par
Số bi đỏ: $20:\dfrac56=24$ (viên).\par
\emph{Thử lại:} lúc đầu $24$ đỏ, $16$ xanh ($24:40=60\%$); lúc sau $24$ đỏ, $36$ xanh ($24:60=40\%$).
\dapso{$24$ viên bi đỏ.}}
\gap
\begin{probbox}{40}{\stars{1}}
Ba vòi cùng chảy đầy một bể. Vòi A chảy được $40\%$ bể, vòi B chảy được bằng $75\%$ lượng nước vòi A chảy, vòi C chảy phần còn lại. Vòi C chảy ít hơn vòi A $54$ lít. Hỏi bể chứa được bao nhiêu lít nước?
\end{probbox}
\probsol{
Vòi B chảy được: $40\%\times75\%=0{,}4\times0{,}75=0{,}3=30\%$ (bể).\par
Vòi C chảy được: $100\%-40\%-30\%=30\%$ (bể).\par
Vòi A chảy nhiều hơn vòi C: $40\%-30\%=10\%$ (bể), ứng với $54$ lít.\par
Bể chứa được: $54:10\%=54:0{,}1=540$ (lít).
\dapso{$540$ lít.}}
\gap
\begin{probbox}{41}{\stars{1}}
Một cửa hàng ngày thứ nhất bán được $40\%$ số gạo, ngày thứ hai bán được $50\%$ số gạo còn lại thì còn $48$ kg. Hỏi lúc đầu cửa hàng có bao nhiêu ki-lô-gam gạo?
\end{probbox}
\probsol{
\emph{Tính ngược từ cuối.} Sau ngày thứ hai còn lại $100\%-50\%=50\%$ số gạo còn sau ngày thứ nhất, ứng với $48$ kg.\par
Số gạo còn sau ngày thứ nhất: $48:50\%=96$ (kg).\par
Số đó bằng $100\%-40\%=60\%$ số gạo lúc đầu, nên lúc đầu có: $96:60\%=96:0{,}6=160$ (kg).
\dapso{$160$ kg gạo.}}
\gap
\begin{probbox}{42}{\stars{1}}
Tìm $x$, biết: $\ x\times125\%+x:4+x\times0{,}5=3{,}4$.
\end{probbox}
\probsol{
$x:4=x\times0{,}25$, nên vế trái bằng $x\times(1{,}25+0{,}25+0{,}5)=x\times2$.\par
$x\times2=3{,}4$, \ nên $x=1{,}7$.
\dapso{$x=1{,}7$.}}
\gap
\begin{probbox}{43}{\stars{1}}
Năm nay tuổi con bằng $25\%$ tuổi mẹ. Sau $6$ năm nữa, tuổi con bằng $\dfrac13$ tuổi mẹ. Tính tuổi của mỗi người hiện nay.
\end{probbox}
\probsol{
Hiện nay: tuổi con $1$ phần, tuổi mẹ $4$ phần; hiệu là $3$ phần.\par
Sau $6$ năm: tuổi con $1$ phần mới, tuổi mẹ $3$ phần mới; hiệu là $2$ phần mới.\par
Hiệu số tuổi không đổi, nên $3$ phần $=2$ phần mới, tức là $1$ phần mới $=1{,}5$ phần.\par
Tuổi con tăng từ $1$ phần lên $1{,}5$ phần: $0{,}5$ phần ứng với $6$ năm, nên $1$ phần là $12$ tuổi.\par
Tuổi con: $12$ tuổi; \ tuổi mẹ: $12\times4=48$ (tuổi).\par
\emph{Thử lại:} sau $6$ năm con $18$ tuổi, mẹ $54$ tuổi, và $18=\dfrac13$ của $54$.
\dapso{con $12$ tuổi, mẹ $48$ tuổi.}}
\gap
\begin{probbox}{44}{\stars{1}}
Một bình chứa $200$ g nước muối, trong đó muối chiếm $5\%$. Cần đổ thêm bao nhiêu gam nước lọc để muối chỉ chiếm $4\%$?
\end{probbox}
\probsol{
Lượng muối trong bình: $200\times5\%=10$ (g). Đổ thêm nước thì lượng muối không đổi.\par
Sau khi đổ thêm, $10$ g muối bằng $4\%$ lượng nước muối: $\ 10:4\%=10:0{,}04=250$ (g).\par
Cần đổ thêm: $250-200=50$ (g nước).
\dapso{$50$ g nước.}}
\gap
\begin{probbox}{45}{\stars{1}}
Số A lớn hơn số B là $25\%$ (của số B). Hỏi số B bé hơn số A bao nhiêu phần trăm (của số A)?
\end{probbox}
\probsol{
Số A bằng $125\%$ số B, tức là $A=\dfrac54\ B$: số B là $4$ phần, số A là $5$ phần.\par
B bé hơn A $1$ phần, mà $1:5=0{,}2=20\%$ (số A).
\dapso{$20\%$.}}
\gap
\begin{probbox}{46}{\stars{1}}
Một lớp học có $45\%$ số học sinh là nam, số học sinh nữ nhiều hơn số học sinh nam $6$ bạn. Hỏi lớp học đó có bao nhiêu học sinh?
\end{probbox}
\probsol{
Số học sinh nữ chiếm: $100\%-45\%=55\%$ (số học sinh cả lớp).\par
Số nữ hơn số nam: $55\%-45\%=10\%$ (số học sinh cả lớp), ứng với $6$ bạn.\par
Số học sinh cả lớp: $6:10\%=60$ (học sinh).
\dapso{$60$ học sinh.}}
\gap
\begin{probbox}{47}{\stars{2}}
Nấm tươi chứa $92\%$ là nước. Sau khi phơi khô, nấm chỉ còn chứa $20\%$ là nước. Hỏi từ $100$ kg nấm tươi thu được bao nhiêu ki-lô-gam nấm khô?
\end{probbox}
\probsol{
Trong $100$ kg nấm tươi, phần ``cái'' (không phải nước) chiếm $100\%-92\%=8\%$: $\ 100\times8\%=8$ (kg).\par
Khi phơi chỉ có nước bay hơi, nên nấm khô vẫn chứa $8$ kg phần ``cái''.\par
Trong nấm khô, phần ``cái'' chiếm $100\%-20\%=80\%$.\par
Lượng nấm khô: $8:80\%=8:0{,}8=10$ (kg).
\dapso{$10$ kg nấm khô.}}
\gap
\begin{probbox}{48}{\stars{1}}
Tổng của hai số là $168$; $40\%$ số thứ nhất bằng $30\%$ số thứ hai. Tìm hai số đó.
\end{probbox}
\probsol{
$\dfrac4{10}$ số thứ nhất $=\dfrac3{10}$ số thứ hai, nên $4$ lần số thứ nhất bằng $3$ lần số thứ hai.\par
Do đó số thứ nhất là $3$ phần, số thứ hai là $4$ phần.\par
Giá trị một phần: $168:7=24$. \ Hai số là: $24\times3=72$ và $24\times4=96$.\par
\emph{Thử lại:} $40\%\times72=28{,}8=30\%\times96$.
\dapso{$72$ và $96$.}}
\gap
\begin{probbox}{49}{\stars{2}}
Một số giảm đi $25\%$. Hỏi phải tăng số mới thêm bao nhiêu phần trăm để được số ban đầu?
\end{probbox}
\probsol{
Coi số ban đầu là $4$ phần. Giảm đi $25\%=\dfrac14$ thì còn $3$ phần.\par
Muốn trở lại số ban đầu phải tăng thêm $1$ phần, tức là bằng $\dfrac13$ số mới.\par
$\dfrac13=1:3=33\dfrac13\%$.
\dapso{$33\dfrac13\%$.}}
\gap
\begin{probbox}{50}{\stars{1}}
An, Bình và Chi góp tiền mua quà. An góp $30\%$ số tiền; Bình góp $40\%$ số tiền còn lại; Chi góp phần còn lại. Chi góp nhiều hơn Bình $14\,000$ đồng. Hỏi mỗi bạn góp bao nhiêu tiền?
\end{probbox}
\probsol{
Sau phần của An còn lại: $100\%-30\%=70\%$ (số tiền).\par
Bình góp: $70\%\times40\%=28\%$ (số tiền); \ Chi góp: $70\%-28\%=42\%$ (số tiền).\par
Chi góp nhiều hơn Bình: $42\%-28\%=14\%$ (số tiền), ứng với $14\,000$ đồng.\par
Số tiền mua quà: $14\,000:14\%=14\,000:0{,}14=100\,000$ (đồng).\par
An góp $30\,000$ đồng; \ Bình góp $28\,000$ đồng; \ Chi góp $42\,000$ đồng.
\dapso{An $30\,000$ đồng, Bình $28\,000$ đồng, Chi $42\,000$ đồng.}}
